% !TEX root = ../main.tex
% !TEX program = lualatex
\documentclass[../main.tex]{subfiles}
\IfSubfilesClassLoaded{\externaldocument{\subfix{../build/main}}}{}
% =============================================================================
% File: 19_Fleet_Maintenance_Assets.tex
% =============================================================================
\begin{document}
\IfSubfilesClassLoaded{\chapter{EDO Study Guide}}{}
\section{Fleet Maintenance Assets}

\subsection{Summary}
\begin{description}
  \item[\textbf{Levels.}] Fleet maintenance is organized into organizational (O), intermediate (I), and depot (D) levels, each matched to increasing complexity, equipment, and workforce specialization~\autocite{edo-4-1-1-fleet-maint-assets-2025}.
  \item[\textbf{Activities.}] O-level is ship's force; I-level includes tenders and shore intermediate facilities; D-level includes public and private shipyards, regional maintenance centers, and ship repair facilities~\autocite{edo-4-1-1-fleet-maint-assets-2025}.
  \item[\textbf{SUPSHIP and RMCs.}] SUPSHIPs administer new construction and modernization contracts, while RMCs manage depot projects and provide I-level support and training~\autocite{edo-4-1-1-fleet-maint-assets-2025}.
  \item[\textbf{Chains and funding.}] Reporting chains and funding streams differ across organizations and must be aligned for effective maintenance execution~\autocite{edo-4-1-1-fleet-maint-assets-2025}.
  \item[\textbf{Workforce.}] Maintenance activities rely on blended military, civilian, and contractor workforces with structured training pipelines and significant early-career demographics at shipyards~\autocite{edo-4-1-1-fleet-maint-assets-2025}.
\end{description}

\subsection{Organizational (O) Level Maintenance}
\begin{description}
  \item[\textbf{Definition.}] Planned maintenance from the \ac{pmsplan} and limited repair performed by ship's force~\autocite{edo-4-1-1-fleet-maint-assets-2025}.
  \item[\textbf{Typical work.}] Tasks range from simple lubrication to modular change-outs and valve overhauls~\autocite{edo-4-1-1-fleet-maint-assets-2025}.
  \item[\textbf{Purpose.}] Leverages operator experience to keep the ship as self-sufficient as possible~\autocite{edo-4-1-1-fleet-maint-assets-2025}.
\end{description}

\subsection{Intermediate (I) Level Maintenance}
\begin{description}
  \item[\textbf{Definition.}] Planned and corrective maintenance on shipboard hull, mechanical, and electrical systems and components~\autocite{edo-4-1-1-fleet-maint-assets-2025}.
  \item[\textbf{Typical work.}] Routine calibration, gas turbine change-outs, pump and engine overhauls, corrosion control, and rigging/weight testing~\autocite{edo-4-1-1-fleet-maint-assets-2025}.
  \item[\textbf{Workforce.}] Sailors and civilians at \ac{ima}/\ac{imf} sites gain \ac{nec} qualifications through OJT-driven repair work~\autocite{edo-4-1-1-fleet-maint-assets-2025}.
\end{description}

\subsection{Depot (D) Level Maintenance}
\begin{description}
  \item[\textbf{Definition.}] Major industrial work requiring specialized equipment, machinery, and skill sets~\autocite{edo-4-1-1-fleet-maint-assets-2025}.
  \item[\textbf{Typical work.}] Docking and underwater hull work, shaft repairs, engine rebuilds, modernization and ship alteration installation, and array resurfacing~\autocite{edo-4-1-1-fleet-maint-assets-2025}.
  \item[\textbf{Workforce.}] Executes with skilled civilian and private-sector workforces beyond I-level capability or capacity~\autocite{edo-4-1-1-fleet-maint-assets-2025}.
\end{description}

\subsection{Maintenance Trends}
\begin{itemize}
  \item O-level tasking shifted ashore for some platforms (e.g., \ac{lcs} and DDG-1000), with crews regaining some checks over time~\autocite{edo-4-1-1-fleet-maint-assets-2025}.
  \item Regional consolidation of I- and D-level work led to integrated shipyard and intermediate facilities at Puget Sound and Pearl Harbor~\autocite{edo-4-1-1-fleet-maint-assets-2025}.
  \item A 2020 GAO review highlighted workforce capacity, unplanned work, planning adherence, facility condition, and IT constraints in surface maintenance~\autocite{edo-4-1-1-fleet-maint-assets-2025}.
\end{itemize}

\subsection{Maintenance Organizations}
\begin{description}
  \item[\textbf{O-level.}] Ship's force~\autocite{edo-4-1-1-fleet-maint-assets-2025}.
  \item[\textbf{I-level.}] Submarine tenders, \ac{sfima}, \ac{nssf}, and \ac{trf}~\autocite{edo-4-1-1-fleet-maint-assets-2025}.
  \item[\textbf{D/I-level.}] Naval shipyards and \acp{imf}, U.S. Naval Ship Repair Facility, \acp{rmc}, \ac{supship} (minor repair role), private shipyards, and some warfare center divisions~\autocite{edo-4-1-1-fleet-maint-assets-2025}.
\end{description}

\subsection{Afloat Maintenance Activities}
\begin{description}
  \item[\textbf{Submarine tenders.}] Provide deployable I-level maintenance to submarines, conduct casualty repair, and support Indo-Asia-Pacific sustainment (e.g., USS Emory S. Land and USS Frank Cable in Guam)~\autocite{edo-4-1-1-fleet-maint-assets-2025}.
  \item[\textbf{SFIMA.}] Afloat \ac{ima} capability using carrier and amphibious platforms to maximize battle force self-sustainment while deployed~\autocite{edo-4-1-1-fleet-maint-assets-2025}.
\end{description}

\subsection{Tender Repair Organization (Typical)}
\begin{description}
  \item[\textbf{Repair officer.}] Leads the tender repair organization; an \ac{edo} O-5 billet in many organizations~\autocite{edo-4-1-1-fleet-maint-assets-2025}.
  \item[\textbf{Divisions.}] Hull (R-1), Mechanical (R-2), Electrical (R-3), Electronic (R-4), Radiological Controls (R-5), Repair Services (R-6), Planning (R-7), Outside Machinery (R-8), Quality Assurance (R-9), and Tender Maintenance (R-T)~\autocite{edo-4-1-1-fleet-maint-assets-2025}.
  \item[\textbf{Support.}] Repair administration, ship superintendent, and technical library functions align under repair leadership~\autocite{edo-4-1-1-fleet-maint-assets-2025}.
\end{description}

\subsection{Submarine Support Facilities}
\begin{description}
  \item[\textbf{NSSF.}] Provides I-level maintenance to fast-attack submarines and support craft; works closely with \ac{rsg} and \ac{nrmd}; based in Groton, CT~\autocite{edo-4-1-1-fleet-maint-assets-2025}.
  \item[\textbf{TRF.}] Provides industrial support for incremental overhaul and refits of ballistic- and guided-missile submarines; conducts I-level maintenance during refit cycles in Kings Bay, GA and Bangor, WA~\autocite{edo-4-1-1-fleet-maint-assets-2025}.
  \item[\textbf{TRIPER.}] \ac{trf} supports equipment replacement using a rotatable pool concept under \ac{triper}~\autocite{edo-4-1-1-fleet-maint-assets-2025}.
\end{description}

\subsection{Naval Shipyards}
\begin{description}
  \item[\textbf{Mission.}] Provide affordable, timely, and quality depot-level maintenance, modernization, inactivation, disposal, and emergency repair of carriers and submarines~\autocite{edo-4-1-1-fleet-maint-assets-2025}.
  \item[\textbf{Capabilities.}] Maintain dry docks, piers, and production shops; sustain a skilled workforce; provide organic nuclear and conventional repair capacity~\autocite{edo-4-1-1-fleet-maint-assets-2025}.
  \item[\textbf{Product lines.}] Submarine availabilities/defueling/refueling; \ac{cvn} availabilities; and combined inactivation, reactor compartment disposal, and hull recycling at specific yards~\autocite{edo-4-1-1-fleet-maint-assets-2025}.
  \item[\textbf{Locations and Specializations.}] The Navy operates four public Naval Shipyards:
    \begin{itemize}
      \item \textbf{Portsmouth Naval Shipyard (PNSY)} (Kittery, ME): Specializes in Los Angeles- and Virginia-class \acs{ssn} overhaul, repair, and modernization~\autocite{navsea-pnsy-about-2026}.
      \item \textbf{Norfolk Naval Shipyard (NNSY)} (Portsmouth, VA): Specializes in aircraft carrier (\ac{cvn}) and submarine (\ac{ssn}/\ac{ssbn}) maintenance, modernization, and inactivation.
      \item \textbf{Puget Sound Naval Shipyard \& IMF (PSNS \& IMF)} (Bremerton, WA): Specializes in \ac{cvn} and submarine (\ac{ssn}/\ac{ssbn}/\ac{ssgn}) maintenance, defueling, inactivation, reactor compartment disposal, and hull recycling. It is the only shipyard that recycles nuclear-powered ships.
      \item \textbf{Pearl Harbor Naval Shipyard \& IMF (PHNSY \& IMF)} (Pearl Harbor, HI): Specializes in submarine (\ac{ssn}) maintenance and modernization, and emergent aircraft carrier (\ac{cvn}) support for the Pacific Fleet.
    \end{itemize}
\end{description}

\subsection{Ship Repair Facility (Japan)}
\begin{description}
  \item[\textbf{SRF/JRMC.}] A non-nuclear depot facility that executes I- and D-level availabilities and provides fleet technical assistance and assessments for Seventh Fleet support~\autocite{edo-4-1-1-fleet-maint-assets-2025}.
  \item[\textbf{Services.}] Planning/execution of availabilities by organic workforce (Yokosuka) and contract oversight (Sasebo), plus dive locker and underwater husbandry services~\autocite{edo-4-1-1-fleet-maint-assets-2025}.
  \item[\textbf{Product lines.}] Homeported ship support in Yokosuka and Sasebo, plus emergent repair brokered by \ac{tycom}~\autocite{edo-4-1-1-fleet-maint-assets-2025}.
\end{description}

\subsection{Supervisor of Shipbuilding (SUPSHIP)}
\begin{description}
  \item[\textbf{Primary mission.}] Administer Navy new construction, nuclear repair, and modernization contracts~\autocite{edo-4-1-1-fleet-maint-assets-2025}.
  \item[\textbf{Key roles.}] Acts as \ac{aco} for shipbuilding/conversion contracts, manages schedules, conducts technical oversight, administers funds, and coordinates logistics support~\autocite{edo-4-1-1-fleet-maint-assets-2025}.
  \item[\textbf{RCOH.}] Conducts \ac{cvn} \ac{rcoh} at Newport News Shipbuilding and is expanding repair responsibilities at nuclear SUPSHIPs~\autocite{edo-4-1-1-fleet-maint-assets-2025}.
\end{description}

\subsection{Major Platform Private Shipyards}
\begin{itemize}
  \item Huntington Ingalls Industries (Newport News Shipbuilding, Ingalls Shipbuilding)~\autocite{edo-4-1-1-fleet-maint-assets-2025}.
  \item General Dynamics (NASSCO, Bath Iron Works, Electric Boat)~\autocite{edo-4-1-1-fleet-maint-assets-2025}.
  \item Fincantieri Marinette Marine, Austal USA~\autocite{edo-4-1-1-fleet-maint-assets-2025}.
\end{itemize}

\subsection{Regional Maintenance Centers}
\begin{description}
  \item[\textbf{Mission.}] Provide D-level project/contract management, I-level repairs and training, fleet technical assistance, \ac{tsra}, underwater ship husbandry, regional calibration, and logistics support~\autocite{edo-4-1-1-fleet-maint-assets-2025}.
  \item[\textbf{Organization.}] RMCs organize around contracts, logistics, financial, corporate operations, environmental safety, engineering and planning, waterfront operations, quality assurance, and production resources~\autocite{edo-4-1-1-fleet-maint-assets-2025}.
  \item[\textbf{Locations.}] Commander, Navy Regional Maintenance Center (\ac{cnrmc}) oversees six \acp{rmc}:
    \begin{itemize}
      \item \textbf{Mid-Atlantic Regional Maintenance Center (MARMC)} (Norfolk, VA).
      \item \textbf{Southeast Regional Maintenance Center (SERMC)} (Mayport, FL).
      \item \textbf{Southwest Regional Maintenance Center (SWRMC)} (San Diego, CA).
      \item \textbf{Forward Deployed Regional Maintenance Center (FDRMC)} (Naples, Italy, with detachments in Rota, Spain and Manama, Bahrain).
      \item \textbf{Northwest Regional Maintenance Center (NWRMC)} (Bremerton, WA) -- co-located and integrated with PSNS \& IMF.
      \item \textbf{Hawaii Regional Maintenance Center (HRMC)} (Pearl Harbor, HI) -- co-located and integrated with PHNSY \& IMF.
    \end{itemize}
  \item[\textbf{Cooperative Support.}] Public Naval Shipyards are government-owned and focused on nuclear-powered vessels (submarines and carriers). RMCs manage conventional surface ship maintenance and modernization, executing depot repairs primarily through private contractors (using Multi-Ship, Multi-Option contracts) while providing intermediate-level (I-level) military repair capacity and fleet technical assistance.
  \item[\textbf{EDO Support roles.}] \acp{edo} provide technical, business, and leadership expertise at these activities:
    \begin{itemize}
      \item At shipyards, \acp{edo} serve as Shipyard Commanders (COs), Production Officers, Project Superintendents, Quality Managers, and Chief Test Engineers.
      \item At RMCs, \acp{edo} serve as Commanding Officers, Waterfront Chief Engineers (WFCHENGs), Project Managers, and Contracting Officer's Representatives.
      \item At SUPSHIPs, \acp{edo} serve as Supervisors, Deputy Supervisors, and Administrative Contracting Officers (ACOs) administering new construction and repair contracts at private yards.
    \end{itemize}
\end{description}

\subsection{Funding and Reporting Chains}
\begin{description}
  \item[\textbf{Modernization funds.}] Maintenance and modernization use separate funding streams, including D-Alt and K-Alt modernization accounts; an \ac{opn} pilot supports some private contracted availabilities~\autocite{edo-4-1-1-fleet-maint-assets-2025}.
  \item[\textbf{Operational chain.}] \ac{cno} delegates readiness through fleet commanders and \acp{tycom} who authorize work and manage maintenance budgets~\autocite{edo-4-1-1-fleet-maint-assets-2025}.
  \item[\textbf{Technical chain.}] \ac{navsea} provides engineering standards and technical authority for maintenance execution at \acp{rmc} and \acp{nsy}~\autocite{edo-4-1-1-fleet-maint-assets-2025}.
\end{description}

\subsection{Responsibilities by Organization}
\begin{description}
  \item[\textbf{OPNAV N83B.}] Establishes ship maintenance policy and procedures, allocates shipyard man-days, and publishes depot availability schedules~\autocite{edo-4-1-1-fleet-maint-assets-2025}.
  \item[\textbf{Fleet commanders.}] Assess material readiness and resource maintenance funding through Fleet Maintenance Officers~\autocite{edo-4-1-1-fleet-maint-assets-2025}.
  \item[\textbf{TYCOM.}] Acts as fleet commander agent for assigned platforms; authorizes work and screens maintenance requirements to \acp{rmc} and depots~\autocite{edo-4-1-1-fleet-maint-assets-2025}.
  \item[\textbf{NAVSEA 04/05.}] Develop technical requirements and engineering standards, approve maintenance strategies, and set industrial activity policy~\autocite{edo-4-1-1-fleet-maint-assets-2025}.
  \item[\textbf{SEA 21 and CNRMC.}] Integrate maintenance strategies, modernization plans, training, and in-service life-cycle management for surface ships~\autocite{edo-4-1-1-fleet-maint-assets-2025}.
\end{description}

\subsection{Workforce Composition and Training}
\begin{description}
  \item[\textbf{Workforce mix.}] Maintenance activities employ military, civilian, and contractor workforces, with \ac{ima} and \ac{imf} pipelines providing OJT-based skill development~\autocite{edo-4-1-1-fleet-maint-assets-2025}.
  \item[\textbf{Shipyard demographics.}] FY17 Q1 data show a large early-career workforce: Portsmouth 40.7\%, Norfolk 46.2\%, Puget Sound 43.3\%, and Pearl Harbor 35.8\% with five years or less experience~\autocite{edo-4-1-1-fleet-maint-assets-2025}.
\end{description}

\subsection{Quick Review}
\begin{itemize}
  \item Maintenance levels are O, I, and D; each level increases in complexity and industrial capability~\autocite{edo-4-1-1-fleet-maint-assets-2025}.
  \item I-level activities include tenders, \ac{sfima}, \ac{nssf}, and \ac{trf}; D-level includes shipyards, \ac{srf}, \acp{rmc}, and private yards~\autocite{edo-4-1-1-fleet-maint-assets-2025}.
  \item SUPSHIPs administer new construction and modernization contracts and hold a minor repair role~\autocite{edo-4-1-1-fleet-maint-assets-2025}.
\end{itemize}

\IfSubfilesClassLoaded{
  \RenewDocumentCommand{\entryname}{}{\textbf{\color{Modern} Acronym}}
  \RenewDocumentCommand{\descriptionname}{}{\textbf{\color{Modern} Definition}}
  \printnoidxglossary[
    type=\acronymtype,
    title=Acronyms,
    style=long-booktabs
  ]}{}
\end{document}
